\section{Continuum-to-Discrete Self-Consistency}\label{sec:conversion}%
% These labels carried subsections until this section was written as continuous
% prose. They are cross-referenced from five other sections and are kept on the
% section itself rather than renamed there.
\label{sec:coarsening}\label{sec:why-discrete}\label{sec:transfer}%
\label{sec:gb-loops}\label{sec:ledger}

This section is sense~(b) of \cref{sec:intro}: where sense~(a) is a statement
\emph{inside} the continuum model, this one is \emph{between} two descriptions
of a microstructure, that the map from continuum fields to discrete loops
preserves what makes them the same microstructure. The two are independent --- a
model may generate every efficiency from its diffusion tensors and still convert
to a population storing the wrong number of defects.

The second is needed because every cluster-dynamics formulation is a mean-field
theory of a population of discrete objects. Its kernels assume absorbers placed
independently and at random, so that overlap depends on the volume fraction
alone --- a dilute-limit statement known to degrade as the swept fraction
becomes
appreciable~\citep{BrailsfordBullough1972,Mansur1978,GolubovSingh2001}. Three
mechanisms take it apart, at different doses. Coalescence destroys spatial
independence, a merged loop occupying the union of two former positions, so the
Avrami estimate stops estimating anything. Local elastic interaction breaks the
far-field assumption behind \cref{eq:Zsk}, two loops within a few radii sharing
a stress field whose drift term is not the sum of the two isolated
ones~\citep{marchrico2023,marchrico2025}. And gradients break the premise that
the diffusion field has relaxed between absorbers --- the failure
\cref{sec:intro} describes, which spatial resolution repairs for the mobile
species but not for the loops.

The model detects the first itself, which is the useful property: $\varphi_{LL}$
of \cref{eq:avrami} is its own estimate of the probability that a loop has met
another, and approaching unity it is of order its own maximum and carries no
information. The coarsening dose
\begin{equation}
  d_{\mathrm{coarsen}}
  = \min\bigl\{\gamma:\ \varphi_{LL}(\gamma)\ge\varphi^{\star}
    \ \text{on a fraction}\ \ge f^{\star}\ \text{of interior nodes}\bigr\}
  \label{eq:dcoarsen}
\end{equation}
is therefore read off one of the model's own state variables rather than imposed
from outside, at no cost. Which family reaches it first is a property of the
formulation and not the material: between calculations differing only in whether
the efficiencies come from the diffusion tensors or from fitted constants, the
family that coarsens first swaps.

Repairing the dilute limit is not the only reason to want a discrete
representation; it is also the interface to micromechanics. Dislocation dynamics
computes the stress a population exerts and its resistance to a gliding line, a
channeling model needs which obstacles a channel has swept, a fracture model the
local obstacle spacing --- and none takes a density as input, but positions,
radii, Burgers vectors and habit planes. The objective is therefore stronger
than a one-way export: that the two representations be equivalent within stated
tolerances, so a calculation may run in either, or in a hybrid, with the choice
made on dose and cost rather than physics. That is what attaches conservation
laws to the conversion instead of making it a rescaling.

Those requirements cannot all be met exactly, and their ranking is a modeling
decision:

\begin{enumerate}[itemsep=2pt]
\item Stored defects, exactly: $\sum_\ell m_\ell$ must equal the continuum
      content of the region the loops replace, a transfer that loses defects
      appearing downstream as a spurious sink.
\item Loop number, to the integer draw: a continuum density gives a real
      expected count and a discrete population an integer one, so the best
      available is $|N_{\mathrm{draw}}-\Lambda|\le\tfrac12$.
\item The dispersion $\Delta_k$, to the truncation of the size sampling --- a
      requirement that did not exist before \cref{sec:moments}, a monodisperse
      family having nothing to reproduce.
\item Sink strength, reported but not asserted on: the line length a discrete
      population presents differs from $\rho^k_{sL}$ for physical rather than
      numerical reasons, ellipticity among them.
\end{enumerate}

Only the eight loop families convert. The stacking-fault pyramid has no Burgers
vector and no line, hence no discrete representation at all; folding its content
into $c_f$ would report basal vacancies as dislocation content the solver was
then free to move.

The conversion proceeds in four steps. The fields are known at unevenly spaced
nodes, so each must first be given the volume it stands for, which a restricted
Voronoi tessellation supplies --- restricted in that sample points outside the
crystal are rejected rather than charged to the nearest boundary node. On a hexagonal prism the bounding box is $4/3$ of the crystal, so
using it would over-count by a third and scatter loops through six empty wedges.

The tessellation is never built explicitly. $V_j$ is estimated by Monte Carlo:
points are drawn uniformly in the bounding box of the node cloud, those failing
the half-space test $\bm{N}\xv\le\bm{b}$ of the convex hull are discarded, and each
survivor is assigned to its nearest node by a $k$-d tree search, so that
$V_j=(M_j/M)\,V_\Omega$ with $M_j$ the count at node $j$ and $M$ the number of
accepted samples. The estimator is unbiased and converges as $M^{-1/2}$; the
production runs use $M=4\times10^{6}$, which on the $9.1\times10^{4}$ nodes here
is some 44 samples per node on average and a few percent on a typical cell.
Nearest-node assignment is what makes the cells a genuine Voronoi partition,
and the only condition the construction has to meet is that they tile the
crystal exactly --- which they do by construction, every accepted sample being
charged to exactly one node, so that $\sum_j V_j=V_\Omega$ identically and no
volume is created or lost however coarse the sampling is. Sampling error moves
volume \emph{between} neighbouring cells; it cannot change their sum.

Two properties of the partition are worth stating because they are easy to
assume wrongly. First, the cells are \emph{not} the finite elements. The
tessellation is built on the cluster-dynamics nodes, which are the trial-function
nodes of the concentration field and include the mid-edge nodes of every
second-order element; the elements are tetrahedra, the Voronoi cells are convex
polyhedra around points, and neither is a refinement of the other. What the two
share is only the point set. Nodal volumes are not read from the mesh because
the output the conversion consumes carries node positions and no connectivity,
and a Voronoi volume is the natural quantity to reconstruct from positions
alone. Second, a node can legitimately receive zero volume --- on the runs here
about $12\%$ do, being nodes whose neighbours crowd them out at a refinement
boundary. That is the correct answer rather than a failure: a node with no
surrounding volume contributes nothing to a volume average, and it is given
weight zero.

The expected number of loops of family $(s,k)$ at node $j$ is then
\begin{equation}
  \lambda^k_{sj}=\frac{n^k_{sL}(\xv_j)}{\Omega}\,V_j,
  \qquad
  \Lambda^k_s=\sum_j\lambda^k_{sj},
  \qquad
  N_{\mathrm{draw}}=\operatorname{round}\bigl(\Lambda^k_s\bigr).
  \label{eq:counts}
\end{equation}
Both factors are needed: sampling on $n^k_{sL}$ alone would over-populate a
refined region in exact proportion to how well it was resolved. Rounding rather
than drawing a Poisson total is deliberate --- $\operatorname{Var}=\Lambda$
there, and a dose sweep whose count fluctuates by $\sqrt\Lambda$ between
snapshots is dominated by counting noise.

Almost no cell expects a whole loop, and the construction does not require one.
At 10\,dpa the basal family converts to some $7\times10^{2}$ loops spread over
$9\times10^{4}$ nodes, so the typical $\lambda^k_{sj}$ is of order $10^{-2}$ and
every cell is empty in expectation. That is why the draw is taken once for the
family and not once per node: the rounded total $N_{\mathrm{draw}}$ fixes how
many loops exist, and the nodes are then sampled with replacement from the
multinomial with probabilities $\lambda^k_{sj}/\Lambda^k_s$, so a cell with
$\lambda^k_{sj}=10^{-2}$ simply wins a loop one time in a hundred rather than
contributing a hundredth of one. The continuum density is recovered in the mean
over that draw, exactly, and the alternative --- rounding node by node --- would
not recover it at all: it would return zero everywhere and convert nothing,
which is the failure mode this ordering exists to avoid. What a small
$\lambda^k_{sj}$ does cost is spatial resolution of the population, since a
family whose total is a few hundred loops cannot resolve structure finer than
the spacing those few hundred points can express, and it is one reason
\cref{sec:results} reads the conversion as a check on the mean field rather than
as a prediction of where individual loops lie.

Each loop's size is drawn from its home node's log-normal
\cref{eq:lognormal} rather than set to the node mean, two moments having left
nothing to draw from and every loop at a node identical. The sample is rescaled
so that $\sum_\ell m_\ell$ matches the continuum content exactly, which is
requirement~1, leaving a residual at machine precision. Geometry closes the
construction: a loop becomes a polygon on its habit plane of an area matching
$\pi R^{2}$ rather than of that circumradius --- an inscribed hexagon has only
$82.7\%$ of the disc, and it is the stored content that must match ---
elliptical with an axis ratio set by the anisotropy, and centered by rejection
sampling inside the crystal.

What the domain refuses must be predicted rather than discovered. A
discrete-dislocation solver silently declines a loop whose nodes leave the
grain, so the conversion predicts that and keeps the refused share in the
continuum --- counted in defects, not loops, since on one measured transfer one
refusal in six was $83\%$ by count and $98\%$ by stored defects. On the geometry
of \cref{sec:results} the basal population is effectively infeasible to transfer
above $\sim0.1\,$dpa, a \cloop{} loop of mean radius $25\,$nm overhanging any
face it lies within $25\,$nm of.

The grain boundary requires two treatments, the objects differing in kind. For
the mobile species it is a Dirichlet condition: \cref{eq:bc} pins the
concentrations at thermal equilibrium, absorption is the flux
$-\bm D^{x}\nabla c^{x}\!\cdot\!\bm n$, and the depleted layer is solved for
rather than parameterized, its width set by the local diffusivity and sink
density and evolving through the transient. A loop cannot be given a
boundary condition, being untransported with no flux to set. What it has is a size, and one of radius $R^k=\lambda_k\sqrt m$ centered a
distance $x$ from the nearest face intersects it when
\begin{equation}
  m\;\ge\;m_{gb}(x)=\left(\frac{x}{\lambda_k}\right)^{2},
  \label{eq:mgb}
\end{equation}
so the absorbed fraction of each moment is the mass of the family's size
distribution above that threshold. With $\Phi^{(j)}_k(m_\theta)$ the fraction of
the $j$-th moment of \cref{eq:lognormal} above $m_\theta$,
\begin{equation}
  \Phi^{(j)}_k(m_\theta)
  =\tfrac12\operatorname{erfc}\!\left(
     \frac{\ln m_\theta-\mu_k-j\,s_k^{2}}{s_k\sqrt2}\right),
  \label{eq:Phij}
\end{equation}
the three removals of \cref{eq:block-n,eq:block-c,eq:block-q} are
\begin{equation}
  \Theta^k_0=\Phi^{(0)}_k\bigl(m_{gb}\bigr)\,n^k_{sL},
  \qquad
  \Theta^k_1=\Phi^{(1)}_k\bigl(m_{gb}\bigr)\,c^k_{sL},
  \qquad
  \Theta^k_2=\Phi^{(2)}_k\bigl(m_{gb}\bigr)\,q^k_{sL},
  \label{eq:gbloop}
\end{equation}
at a rate $\nu_{gb}$ taken large, so the touching fraction is removed
essentially instantaneously and the denuded-zone width is a prediction rather
than a rate parameter.

Three properties of \cref{eq:gbloop} are worth stating. It is automatically
realizable: what remains after removing the mass above a threshold is a genuine
sub-measure, so $\Delta_k\ge1$ holds by construction. The defects do not return to the pool, the whole difference
between it and the dissolution channel that credits freed vacancies back to
$\spvec c_M$: a loop swallowed by a free surface takes its defects out of the
crystal, so \cref{eq:gbloop} is booked to a grain-boundary ledger alone. Adding
it to the dissolution term too would return the atoms twice, a conservation
residual of the wrong sign the only symptom. And the
distribution must be truncated at the crystal, the log-normal having unbounded
support where the crystal has not: uncapped, the far tail of
\cref{eq:lognormal} holds loops larger than the grain, which \cref{eq:mgb}
faithfully absorbs at every distance, reporting absorption at a grain center no
loop can reach. Capping at $m_{\max}=(R_{\max}/\lambda_k)^{2}$ with $R_{\max}$
the inradius replaces \cref{eq:Phij} by the mass between threshold and cap over
the mass below it,
\begin{equation}
  \widetilde\Phi^{(j)}_k(m_\theta)
  =\frac{\Phi^{(j)}_k(m_\theta)-\Phi^{(j)}_k(m_{\max})}
         {1-\Phi^{(j)}_k(m_{\max})},
  \label{eq:trunc}
\end{equation}
which reduces to $\Phi^{(j)}_k$ exactly as $m_{\max}\to\infty$. The carried
moments are still inverted for $(\mu_k,s_k)$ as if untruncated, so closure and
support disagree at order $\Phi^{(j)}_k(m_{\max})$; reconciling them would mean
inverting a truncated log-normal from truncated moments, which has no closed
form. The inconsistency is bounded by the discarded tail, which is the tail with
no business existing. Measured on \cref{sec:results}, truncation leaves the near field
unchanged to four digits, sets the absorption rate $1\,\mu$m from a boundary to
zero where it had been finite, and raises it at $200\,$nm by $4\%$.

Whether this holds together is settled by auditing every transfer on three
totals, each summed over the crystal whichever representation holds them:
\begin{align}
  \text{(I) loop number} &&
    N &= \sum_j n^k_{sj}V_j \;+\; \#\{\text{discrete loops}\}, \nonumber\\
  \text{(II) stored defects} &&
    C &= \sum_j \frac{c^k_{sj}}{\Omega}V_j
         \;+\; \sum_\ell \frac{\pi R_\ell^{2}(\bvec\!\cdot\!\nvec)_k}{\Omega},
         \label{eq:ledger}\\
  \text{(III) sink strength} &&
    S &= \sum_j 2\pi R^k_{sj}n^k_{sj}V_j \;+\; \sum_\ell 2\pi R_\ell . \nonumber
\end{align}
(I) and (II) are conservation statements, holding to the sampling tolerance of
one loop and to machine precision. (III) is a modeling statement, and the one
that catches mistakes: a discontinuity in it means the two representations
disagree about what a loop of a given content is, a crystallography error rather
than a numerical one. The residual sink discontinuity fell from a
factor of $4.85$ to $2\%$ by removing an inconsistency in the \cloop{} Burgers
vector --- the continuum sizing a basal loop with the full $[0001]$ where the
discrete side used $\tfrac12[0001]$ --- and not by tuning.

One constraint follows: a transfer is all-or-nothing per family. The continuum
carries one mean size per family per node, so a size-selective transfer is not
representable. More importantly, a family discrete \emph{and} still present in
the continuum is a sink twice over, once through its Green's function and once
through \cref{eq:KL}, and a loop whose own sink remains in the continuum samples
its own depletion field at its own line. Zeroing the transferred family removes
both, and they are the same fault.
